% 计数与长度的不交相加：(a)五点权重0.5,1,2,0.5,1，A={1,2,3}总权重3.5，B={4,5}总权重1.5；(b)E=[0,2)∪[3,4)的两种分割，共同细分切点0,0.5,1.5,2,3,3.5,4，总长度3。
\begin{tikzpicture}[x=1cm,y=1cm,
  pt/.style={circle,draw=FigureStroke,fill=FigureBlueFill,line width=.7pt,inner sep=0pt},
  grpA/.style={draw=FigureBlue,rounded corners=4pt,line width=.7pt},
  grpB/.style={draw=FigureGreen,rounded corners=4pt,line width=.7pt},
  note/.style={font=\FigureNoteFont\fontsize{7.5}{9}\selectfont,text=FigureMuted},
  seg/.style={line width=3pt},
  tick/.style={draw=FigureGrid,line width=.6pt}]
\begin{scope}
  \node[pt,minimum size=3.2mm] (p1) at (0,0) {};
  \node[pt,minimum size=4.5mm] (p2) at (0.7,0) {};
  \node[pt,minimum size=6.4mm] (p3) at (1.5,0) {};
  \node[pt,minimum size=3.2mm,fill=FigureGreenFill] (p4) at (2.6,0) {};
  \node[pt,minimum size=4.5mm,fill=FigureGreenFill] (p5) at (3.25,0) {};
  \foreach \i/\w in {1/0.5,2/1,3/2,4/0.5,5/1}{\node[font=\fontsize{7.5}{9}\selectfont] at ($(p\i)+(0,-.62)$) {$\w$};\node[font=\fontsize{8}{9}\selectfont] at ($(p\i)+(0,.55)$) {$\i$};}
  \draw[grpA] (-.35,-.95) rectangle (1.95,.85);
  \draw[grpB] (2.3,-.95) rectangle (3.6,.85);
  \node[text=FigureBlue,anchor=south west] at (-.35,.85) {$A$};
  \node[text=FigureGreen,anchor=south west] at (2.3,.85) {$B$};
  \node[note] at (.8,-1.25) {$\#A=3$，权重$3.5$};
  \node[note] at (2.95,-1.25) {$\#B=2$，权重$1.5$};
  \node at (1.6,-1.85) {(a) 计数与加权计数};
\end{scope}
\begin{scope}[xshift=5.15cm]
  \def\u{1.2}
  \foreach \x in {0,0.5,1.5,2,3,3.5,4}{\draw[tick,dashed] (\x*\u,-.55)--(\x*\u,.9);}
  \draw[seg,FigureBlue] (0,.55) -- (1.5*\u,.55);
  \draw[seg,FigureBlueFill!60!FigureBlue] (1.5*\u,.55) -- (2*\u,.55);
  \draw[seg,FigureBlue] (3*\u,.55) -- (4*\u,.55);
  \foreach \x in {0,1.5,2,3,4}{\draw[FigureStroke,line width=.6pt] (\x*\u,.42)--(\x*\u,.68);}
  \node[note,anchor=east] at (-.15,.55) {分割一};
  \draw[seg,FigureGreen] (0,-.15) -- (.5*\u,-.15);
  \draw[seg,FigureGreenFill!40!FigureGreen] (.5*\u,-.15) -- (2*\u,-.15);
  \draw[seg,FigureGreen] (3*\u,-.15) -- (3.5*\u,-.15);
  \draw[seg,FigureGreenFill!40!FigureGreen] (3.5*\u,-.15) -- (4*\u,-.15);
  \foreach \x in {0,0.5,2,3,3.5,4}{\draw[FigureStroke,line width=.6pt] (\x*\u,-.28)--(\x*\u,-.02);}
  \node[note,anchor=east] at (-.15,-.15) {分割二};
  \draw[FigureStroke,line width=.6pt,->] (-.15,-.75)--(4*\u+.3,-.75);
  \foreach \x/\l in {0/0,0.5/0.5,1.5/1.5,2/2,3/3,3.5/3.5,4/4}{\node[font=\fontsize{7}{8}\selectfont,below] at (\x*\u,-.75) {$\l$};}
  \node[note] at (2*\u,1.15) {$1.5+0.5+1=0.5+1.5+0.5+0.5=3$};
  \node at (2*\u,-1.85) {(b) 有限个不交区间的长度};
\end{scope}
\end{tikzpicture}
